\section{Introduction}
\label{sec:introduction}

Offline reinforcement learning (RL) seeks a policy from a fixed dataset without additional environment interaction. Its central difficulty is not merely to select a high-value action, but to determine how far policy improvement can proceed before the learned critic becomes unreliable. TD3+BC regularizes a deterministic actor toward dataset actions while using a Polyak copy of that actor in bootstrapped critic targets \citep{fujimoto2021minimalist}. The same policy branch therefore determines both the action shift entering critic training and the reach of the evaluated policy. Aggressive improvement can move target actions into regions where critic error is amplified, whereas conservative improvement can leave useful policy improvement unrealized.

We study this coupling through \emph{multi-step proximal policy improvement} (MPI), a re-centered refinement architecture for TD3+BC. MPI represents a nominal improvement budget with $K$ sequential actor updates. The first actor is anchored to dataset actions, and each later actor is anchored to the stop-gradient output of its predecessor. A target copy of the first refined actor supplies the next actions in TD targets, while the final actor is used for evaluation. The architecture thus assigns critic-target construction and deployment to distinct policy branches without changing the critic update rule.

Let $D_{\mathrm{critic}}$ denote the squared displacement of the critic-coupled first-hop branch and let $D_{\mathrm{final}}$ denote that of the evaluated final policy. A one-hop actor identifies these quantities; MPI assigns them to $\mu_1$ and $\mu_K$, respectively. This separation can increase final-policy reach at comparable first-hop movement or contract both movements near an unstable training outcome. It also introduces a countervailing mismatch because the critic target follows the $\mu_1$ branch while evaluation follows $\mu_K$. The method therefore exposes a design axis between target-policy displacement and deployment reach rather than guaranteeing that either quantity is independently optimal.

For deterministic actors, the squared action anchor in TD3+BC is a Wasserstein--2 movement cost in the Dirac-policy limit. Under a fixed critic, inactive action projection, accurate subproblem solutions, and small \emph{realized} steps, explicit and implicit updates are first-order discretizations of the same Wasserstein gradient flow. At finite steps, a proximal comparison yields an improvement margin only relative to the learned critic. At large nominal budgets, projection, neural actor fitting, critic drift, saturation, and the seeded optimization trajectory determine the observed end-to-end outcome. We keep these local, finite-step, and coupled regimes distinct throughout the analysis.

We evaluate this formulation with a complete fixed-budget sweep over nine D4RL locomotion tasks, fourteen budgets, and two paired seeds. Three proximal hops reduce large-budget low-return cells from $35/63$ for TD3+BC to $17/63$ and primarily improve the lower tail rather than uniformly shifting scores. Three scale-matched projected linearized hops reduce failures to $24/63$, so the stability effect is not unique to the proximal loss. Direct next-state diagnostics show that three-hop MPI lowers deterministic target-policy displacement in $88/90$ paired comparisons with TD3+BC, while its current-state final-policy displacement proxy is lower in only $36/90$. Low-return runs additionally exhibit orders-of-magnitude increases in TD-error tails, larger realized movements, and critic gains that reverse sign under a controlled common critic. In a paired shadow-critic intervention, final-hop routing worsens simulator-referenced calibration relative to first-hop routing in all eight targeted runs.

Our contributions are as follows.
\begin{itemize}
    \item We formulate MPI for TD3+BC as a re-centered actor composition that assigns target-policy displacement and final-policy reach to distinct branches.
    \item We give a Wasserstein proximal interpretation, derive an inexact finite-step margin with an explicit critic-error term, bound first-hop TD-target perturbation, and separate these statements from the fixed-critic small-step flow limit.
    \item Through complete proximal and scale-matched linearized sweeps, direct target-policy measurements, and paired shadow-critic diagnostics, we show that substepping broadens the empirical large-budget stability envelope and that low-return outcomes coincide with critic miscalibration rather than local truncation error alone.
\end{itemize}
